\documentclass[10pt,a4paper]{amsart}
\usepackage[margin=19mm]{geometry}
\usepackage{amsmath,amssymb,mathtools}
\usepackage[scr=rsfs,cal=euler]{mathalfa}
\usepackage{xfrac}
\usepackage[unicode,hidelinks,bookmarks=false]{hyperref}
\newcommand{\ie}{i.e.\ }
\newcommand{\cf}{cf.\ }
\newcommand{\resp}{respectively\ }
\newcommand{\Subsection}{\subsection}
\providecommand{\nobreakdash}{\nobreak\hbox{-}}
\setlength{\emergencystretch}{2em}
\multlinegap0pt
\usepackage{amssymb}
\usepackage{mathtools}
\usepackage[shortlabels]{enumitem}
\newtheorem{prop}{Proposition}
\newtheorem{lem}{Lemma}
\title{On the problem of large gcd for disjoint residue classes}
\author{Jan Fornal, Yu-Chen Sun}
\date{Source: arXiv:2607.24655v1 (CC BY 4.0)}
\begin{document}
\maketitle

\noindent\emph{Source and licence.} On the problem of large gcd for disjoint residue classes. Version arXiv:2607.24655v1 (CC BY 4.0), distributed by arXiv under the Creative Commons Attribution 4.0 International licence, http://creativecommons.org/licenses/by/4.0/, as recorded in the arXiv licence metadata for this exact version and shown on the abstract page https://arxiv.org/abs/2607.24655v1. Full licence notice: \texttt{licenses/arxiv-2607.24655-CC-BY-4.0.txt}.\par\smallskip

\noindent\textit{Editorial scope.} Counted unit: the article's prop prop:sieve-theoretic-partition (source0.tex lines 324--338). The article's complete original proof of it is delivered with it (source0.tex lines 511--674, one \texttt{proof} environment of 164 lines, which the article places away from the statement). No mathematics is written, repaired or re-derived: the statement and the proof are the article's own lines, in the article's own notation, and no retained line is substituted or re-flowed.  Retained uncounted as the source-local context the proof needs: lines 466--470; lines 472--484.  Nothing was omitted from any retained span, so the package carries no omission record.  No retained line contains a citation, so the wrapper carries no bibliography and no bibliography entry.  No retained line contains a figure environment, an image-inclusion command or drawing code, so no image is delivered and the package declares no asset dependency.  Editorial formatting only, in addition: an AMS \texttt{amsart} preamble that restates the article's own declarations and macros used by the retained lines, namely the environments \texttt{prop}, \texttt{lem} on separate counters, together with the standard packages those lines need.  The retained environments print in this document as Proposition 1, Lemma 1 in order of appearance, whereas the article prints the same items with the numbering of its own full document.  The complete native source of the article as distributed in the arXiv e-print for this exact version is bundled unchanged as \texttt{sources/206131/source0.tex}.
\begin{equation}\label{eq:HjQj}
 H_j=\sum_{p\in\mathcal P_j}\frac1p,
 \qquad
 Q_j=\sum_{p\in\mathcal P_j}\frac1{p^2}.
\end{equation}

\begin{lem}\label{lem:multi-moment}
Suppose that, for every $j\in J$, a real number $z_j$ is chosen with
\begin{equation}
 2\le z_j\le \frac{\mathrm e^{U_j}}2.
\end{equation}
Then
\begin{equation}
 A_I(X,\mathbf t)
 \le X\prod_{j\in J} z_j^{-t_j}
 \exp\!\left(\sum_{j\in J} z_jH_j
              +2\sum_{j\in J}z_j^2Q_j\right).
\end{equation}
\end{lem}

\begin{prop} \label{prop:sieve-theoretic-partition}
  For every sufficiently large integer $d$, there exists a partition:
  \begin{equation}
    [d] = \bigsqcup_{i \in \mathcal{I}} C_i
  \end{equation}
  such that the following two conditions are satisfied:
  \begin{enumerate}
    \item The size of the set $|\mathcal{I}| < \exp(o(\sqrt{\frac{\log d}{\log \log d}}))$.
    \item For each $i \in \mathcal{I}$, one has that:
    \begin{equation}\label{eq:Ti-definition}
      T_i := \frac{1}{d} \max_{m \in C_i} \sum_{e | m} \sum_{\substack{n \in C_i \\ \gcd(n, m) = e}} e
    \end{equation}
    is bounded by $\exp((2 + o(1)) \sqrt{\frac{\log d}{\log \log d}})$.
  \end{enumerate}
\end{prop}

\begin{proof}[Proof of Proposition~\ref{prop:sieve-theoretic-partition}]
Applying Lemma~\ref{lem:multi-moment}, with
$t_j=r_j-\Omega(e_j)$, gives
\begin{align*}
 \left(\prod_{j=1}^{J}e_j\right)
 \#\{c:e_j\mid c_j\text{ for all }j\}
 &\le \frac da\prod_jz_j^{-(r_j-\Omega(e_j))}\\
 &\qquad{}\times
 \exp\!\left(\sum_jz_jH_j+2\sum_jz_j^2Q_j\right).
\end{align*}
Let $\sigma(a)=\sum_{e_0\mid a} e_0$ denote the sum-of-divisors function. Summing over all
small-part divisors $e_0\mid a$ and all $e_j\mid b_j$, we obtain
\begin{align}
&\sum_{e\mid m}e\,\#\{n\in C:e\mid n\}\notag\\
&\le d\frac{\sigma(a)}a
 \exp\!\left(\sum_jz_jH_j+2\sum_jz_j^2Q_j\right)
 \prod_j
 \sum_{e_j\mid b_j}z_j^{-(r_j-\Omega(e_j))}.
\label{eq:class-master-raw}
\end{align}

We simplify the last divisor sum.  Replace $e_j$ by the complementary divisor
$f_j=b_j/e_j$.  Since $\Omega(b_j)=r_j$,
\begin{align*}
 \sum_{e_j\mid b_j}z_j^{-(r_j-\Omega(e_j))}
 &=\sum_{f_j\mid b_j}z_j^{-\Omega(f_j)}\\
 &=\prod_{p^\alpha\parallel b_j}
   (1+z_j^{-1}+\cdots+z_j^{-\alpha})\\
 &\le (1-z_j^{-1})^{-\omega_j(b)}.
\end{align*}
where:
\begin{equation} \label{eq:omega-j-definition}
  \omega_j(b) = \# \{ p \in \mathcal{P}_j : p | b \}
\end{equation}
The small-prime factor is harmless because the set $p\le\mathrm e^{U_0}$ is fixed:
\[
 \frac{\sigma(a)}a
 \le\prod_{p\le\mathrm e^{U_0}}(1-1/p)^{-1} \ll \log \log \log d.
\]
Finally, for $z\ge2$,
\[
 -\log(1-z^{-1})\le z^{-1}+2z^{-2}.
\]
Taking logarithms in \eqref{eq:class-master-raw} yields the central inequality
\begin{equation}\label{eq:central-inequality}
 \log\frac{\mathcal T_i(m)}d
 \le O(1)
 +\sum_j\left(z_jH_j+\frac{\omega_j(b)}{z_j}\right)
 +2\sum_jz_j^2Q_j
 +2\sum_j\frac{\omega_j(b)}{z_j^2}.
\end{equation}
Choose
\begin{equation}
 z_j=\max\left\{2,\sqrt{\frac{\omega_j(b)}{H_j}}\right\}.
\end{equation}
One can easily verify that $z_j$ is at most $\frac{e^{U_j}}{2}$, as we have insisted before.
Therefore:
\begin{equation}\label{eq:main-reduction}
 \sum_j\left(z_jH_j+\frac{\omega_j(b)}{z_j}\right)
 \le 2\sum_j\sqrt{H_j\omega_j(b)}+4\sum_jH_j.
\end{equation}
By the Mertens' second theorem:
\[
 \sum_jH_j\le\sum_{p\le d}\frac1p \ll \log \log d.
\]
Also:
\begin{equation}
  2 \sum_j \frac{\omega_j(b)}{z_j^2} \leq 2 \sum_j H_j
\end{equation}
and:
\begin{equation}
  \begin{aligned}
    \sum_{j} z_j^2 Q_j
    &\leq 4\sum_j Q_j+\sum_j \frac{\omega_j(b) Q_j}{H_j}\\
    &\ll 1+\sum_{j : U_j \leq 2 \log \log d}
      \frac{|\mathcal{P}_j| H_j}{e^{U_j} H_j}
      + \frac{\log d}{(\log d)^2}\\
    &\ll \sum_{j : U_j \leq 2 \log \log d} e^{\eta U_j}\\
    &\ll (\log \log d)^{5/2}
      \exp(2 (\log \log d)^{1/2})
  \end{aligned}
\end{equation}
due to the fact that $\sum_j \omega_j(b) \leq \log d$ and $Q_j \leq e^{-U_j} H_j$ by the definitions of $Q_j, H_j$ and $\omega_j(b)$ (see \eqref{eq:HjQj} and \eqref{eq:omega-j-definition}).
Thus, the remaining task is to estimate $\sum_j\sqrt{H_j\omega_j(b)}$.

Let:
\begin{equation}
 U_*=\log \log d-4\sqrt{\log \log d}.
\end{equation}

If $U_j<U_*$, then
\[
 H_j\le |\mathcal{P}_j|\mathrm e^{-U_j},
 \qquad \omega_j(b)\le |\mathcal{P}_j|.
\]
Therefore:
\begin{equation}\label{eq:low-onebox}
 \sqrt{H_j\omega_j(b)}\le |\mathcal{P}_j|\mathrm e^{-U_j/2}.
\end{equation}
If $p\in\mathcal P_j$, then
\[
 \log p\le U_{j+1}=(1+\eta)U_j,
\]
so
\[
 \mathrm e^{-U_j/2}\le p^{-1/(2(1+\eta))}.
\]
All primes in the low boxes are at most
\[
 X=\exp((1+\eta)U_*).
\]
Using \eqref{eq:low-onebox} and then enlarging a prime sum to an integer sum,
\begin{align*}
 \sum_{U_j<U_*}\sqrt{H_j\omega_j(b)}
 &\le \sum_{p\le X}p^{-1/(2(1+\eta))}\\
 &\ll \frac{1}{\log X} X^{1-\frac{1}{2(1+\eta)}}\\
 &\ll \frac{1}{\log \log d}
 \exp\!\left(\left(\frac12-\eta-4\eta^2\right)\log\log d\right)\\
 &= \frac{(\log d)^{1/2} \exp(- \eta \log \log d - 4 \eta^2 \log \log d) }{\log \log d} \\
 &\leq \frac{(\log d)^{1/2}}{\log \log d}
\end{align*}
For the remaining boxes, Cauchy--Schwarz gives
\begin{align}
 \sum_{U_j\ge U_*}\sqrt{H_j\omega_j(b)}
 &\le
 \left(\sum_{U_j\ge U_*}U_j\omega_j(b)\right)^{1/2}
 \left(\sum_{U_j\ge U_*}\frac{H_j}{U_j}\right)^{1/2}.
\label{eq:CS-main}
\end{align}
We estimate the two factors separately.

Every distinct prime counted by $\omega_j(b)$ is larger than $\mathrm e^{U_j}$. Hence
\begin{equation}\label{eq:log-budget}
 \sum_jU_j\omega_j(b)
 \le\sum_{p\mid b}\log p
 =\log\operatorname{rad}(b)
 \le\log b
 \le \log d.
\end{equation}
If $p\in\mathcal P_j$, then $\log p\le(1+\eta)U_j$, and therefore
\[
 \frac1{U_j}\le\frac{1+\eta}{\log p}.
\]
By the prime number theorem,
\begin{align}
 \sum_{U_j\ge U_*}\frac{H_j}{U_j}
 &\le(1+\eta)\sum_{p>\mathrm e^{U_*}}\frac1{p\log p}\notag\\
 &=\frac{1+o(1)}{U_*}
 =\frac{1+o(1)}{\log \log d}.
\label{eq:tail-budget}
\end{align}
Combining \eqref{eq:log-budget} and \eqref{eq:tail-budget} in
\eqref{eq:CS-main}, we obtain
\begin{equation}
 \sum_{U_j\ge U_*}\sqrt{H_j\omega_j(b)}
 \le(1+o(1))\sqrt{\frac{\log d}{\log \log d}}.
\end{equation}
Combining this with \eqref{eq:central-inequality}, \eqref{eq:main-reduction}, and the estimates above gives, uniformly in $m\in C_i$,
\[
  \log\frac{\mathcal T_i(m)}d
  \leq (2+o(1))\sqrt{\frac{\log d}{\log\log d}}.
\]
Thus $\max_{m\in C_i}\mathcal T_i(m)\leq d\exp((2+o(1))\sqrt{\log d/\log\log d})$, and the desired bound for $T_i$ follows from $dT_i\leq \max_{m\in C_i}\mathcal T_i(m)$.
\end{proof}

\end{document}
